%%% Hlavní soubor. Zde se definují základní parametry a odkazuje se na ostatní části. %%%

%% Verze pro jednostranný tisk:
% Okraje: levý 40mm, pravý 25mm, horní a dolní 25mm
% (ale pozor, LaTeX si sám přidává 1in)
\documentclass[12pt,a4paper]{report}
\setlength\textwidth{145mm}
\setlength\textheight{247mm}
\setlength\oddsidemargin{15mm}
\setlength\evensidemargin{15mm}
\setlength\topmargin{0mm}
\setlength\headsep{0mm}
\setlength\headheight{0mm}
% \openright zařídí, aby následující text začínal na pravé straně knihy
\let\openright=\clearpage

%% Pokud tiskneme oboustranně:
% \documentclass[12pt,a4paper,twoside,openright]{report}
% \setlength\textwidth{145mm}
% \setlength\textheight{247mm}
% \setlength\oddsidemargin{15mm}
% \setlength\evensidemargin{0mm}
% \setlength\topmargin{0mm}
% \setlength\headsep{0mm}
% \setlength\headheight{0mm}
% \let\openright=\cleardoublepage

%% Použité kódování znaků: obvykle latin2, cp1250 nebo utf8:
\usepackage[utf8]{inputenc}

%% Ostatní balíčky
\usepackage{graphicx}
\usepackage{amsthm}

%% Balíček hyperref, kterým jdou vyrábět klikací odkazy v PDF,
%% ale hlavně ho používáme k uložení metadat do PDF (včetně obsahu).
%% POZOR, nezapomeňte vyplnit jméno práce a autora.
\usepackage[ps2pdf,unicode]{hyperref}   % Musí být za všemi ostatními balíčky
\hypersetup{pdftitle=File sharing protocol}
\hypersetup{pdfauthor=Jan Matějek}

%%% Drobné úpravy stylu

% Tato makra přesvědčují mírně ošklivým trikem LaTeX, aby hlavičky kapitol
% sázel příčetněji a nevynechával nad nimi spoustu místa. Směle ignorujte.
\makeatletter
\def\@makechapterhead#1{
  {\parindent \z@ \raggedright \normalfont
   \Huge\bfseries \thechapter. #1
   \par\nobreak
   \vskip 20\p@
}}
\def\@makeschapterhead#1{
  {\parindent \z@ \raggedright \normalfont
   \Huge\bfseries #1
   \par\nobreak
   \vskip 20\p@
}}
\makeatother

% Toto makro definuje kapitolu, která není očíslovaná, ale je uvedena v obsahu.
\def\chapwithtoc#1{
\chapter*{#1}
\addcontentsline{toc}{chapter}{#1}
}

\begin{document}

% Trochu volnější nastavení dělení slov, než je default.
\lefthyphenmin=2
\righthyphenmin=2

%%% Titulní strana práce

\pagestyle{empty}
\begin{center}

\large

Charles University in Prague

\medskip

Faculty of Mathematics and Physics

\vfill

{\bf\Large BACHELOR THESIS}

\vfill

\centerline{\mbox{\includegraphics[width=60mm]{logo.eps}}}

\vfill
\vspace{5mm}

{\LARGE Jan Matějek}

\vspace{15mm}

% Název práce přesně podle zadání
{\LARGE\bfseries File sharing protocol}
% {\LARGE\bfseries A file sharing protocol}

\vfill

% Název katedry nebo ústavu, kde byla práce oficiálně zadána
% (dle Organizační struktury MFF UK)
Department of Applied Mathematics

\vfill

\begin{tabular}{rl}

Supervisor of the bachelor thesis: & Martin Mareš \\
\noalign{\vspace{2mm}}
Study programme: & Computer Science \\
\noalign{\vspace{2mm}}
Specialization: & Programming \\
\end{tabular}

\vfill

% Zde doplňte rok
Prague 2012

\end{center}

\newpage

%%% Následuje vevázaný list -- kopie podepsaného "Zadání bakalářské práce".
%%% Toto zadání NENÍ součástí elektronické verze práce, nescanovat.

%%% Na tomto místě mohou být napsána případná poděkování (vedoucímu práce,
%%% konzultantovi, tomu, kdo zapůjčil software, literaturu apod.)

\openright

\noindent
Dedication.

\newpage

%%% Strana s čestným prohlášením k bakalářské práci

\vglue 0pt plus 1fill

\noindent
I declare that I carried out this bachelor thesis independently, and only with the cited
sources, literature and other professional sources.

\medskip\noindent
I understand that my work relates to the rights and obligations under the Act No.
121/2000 Coll., the Copyright Act, as amended, in particular the fact that the Charles
University in Prague has the right to conclude a license agreement on the use of this
work as a school work pursuant to Section 60 paragraph 1 of the Copyright Act.

\vspace{10mm}

\hbox{\hbox to 0.5\hsize{%
In ........ date ............
\hss}\hbox to 0.5\hsize{%
signature of the author
\hss}}

\vspace{20mm}
\newpage

%%% Povinná informační strana bakalářské práce

\vbox to 0.5\vsize{
\setlength\parindent{0mm}
\setlength\parskip{5mm}

Název práce:
Protokol pro sdílení souborů
% přesně dle zadání

Autor:
Jan Matějek

Katedra:  % Případně Ústav:
Katedra aplikované matematiky
% dle Organizační struktury MFF UK

Vedoucí bakalářské práce:
Mgr. Martin Mareš, Ph.D., Katedra aplikované matematiky
% dle Organizační struktury MFF UK, případně plný název pracoviště mimo MFF UK

Abstrakt:
% abstrakt v rozsahu 80-200 slov; nejedná se však o opis zadání bakalářské práce

Klíčová slova:
% 3 až 5 klíčových slov

\vss}\nobreak\vbox to 0.49\vsize{
\setlength\parindent{0mm}
\setlength\parskip{5mm}

Title:
File sharing protocol
% přesný překlad názvu práce v angličtině

Author:
Jan Matějek

Department:
Department of Applied Mathematics
% dle Organizační struktury MFF UK v angličtině

Supervisor:
Mgr. Martin Mareš, Ph.D., Department of Applied Mathematics
% dle Organizační struktury MFF UK, případně plný název pracoviště
% mimo MFF UK v angličtině

Abstract:
% abstrakt v rozsahu 80-200 slov v angličtině; nejedná se však o překlad
% zadání bakalářské práce

Keywords:
% 3 až 5 klíčových slov v angličtině

\vss}

\newpage

%%% Strana s automaticky generovaným obsahem bakalářské práce. U matematických
%%% prací je přípustné, aby seznam tabulek a zkratek, existují-li, byl umístěn
%%% na začátku práce, místo na jejím konci.

\openright
\pagestyle{plain}
\setcounter{page}{1}
\tableofcontents

%%% Jednotlivé kapitoly práce jsou pro přehlednost uloženy v samostatných souborech
\chapter*{Introduction}
\addcontentsline{toc}{chapter}{Introduction}

Today, in the age of ubiquituous networking, it is as important as ever to be able to share files between
devices. Users own smart cameras and phones, home fileservers, shared document storages etc., and want to
work with data stored on these devices. Moreover, one of the basic requirements of cloud computing, which is
rising in popularity and importance, is the ability of clients to efficiently and securely access their data
that reside on remote servers.

This means that networking protocols on connected devices play a huge role in perceived performance, utility
and feature set of those devices and networks -- especially when we're talking about low-powered devices,
unreliable wireless networks and other non-optimal conditions. However, details and properties of those
protocols are abstracted away from their users. The reason for this is obvious: non-technical users are
interested in features and performance of concrete applications, not in implementation details. But this means
that, for instance, actual low-level security and resilience to attacks remains hidden. Also, less and less
attention is paid to cross-application and cross-platform interoperability.

We decided to examine one of the lower levels in the protocol hierarchy -- basic filesystem-like access to
files and directory structures. We know from experience that existing solutions, while sufficient, are
far from perfect in real-world conditions. Interoperability between different implementations is less than ideal,
and some protocols are better able to cope with degradation in network environment (packet loss, high latency)
than others.

We decided to design a new, modern file transfer protocol, in an attempt to examine possible solutions to
various problems and challenges. We hoped to find out how tradeoffs in the design affect different application
scenarios, how much protocol efficiency depends on usage by applications and how features of a protocol relate
to its vulnerability to attacks. This allowed us to come up with a protocol that is better suited for some
use cases than the protocols available today.

The results of this work are three-fold. First and foremost, it is the protocol design, along with reference
implementations for server and client. Second, the thesis describes in detail the various properties of a file
sharing protocol, available design choices and tradeoffs.  This will be useful for people designing a protocol
best suited for a particular use case. And thirdly, we present a summary evaluation of relevant characteristics
of protocols that are widely in use today.

The work is divided into four chapters. In Chapter 1, we define and explain terms and concepts used throughout
the work. This chapter also examines common issues related to file access protocols, and outlines the various
possible solutions. Chapter 2 presents an overview of popular existing protocols, their properties and
features. Chapter 3 is a specification of the protocol that we designed. It also explains the logic behind our
design decisions, and how these decisions relate to usage scenarios. Finally, Chapter 4 describes the software
that we developed in order to test viability of the new protocol.

%\chapter{Title of the first chapter}

\section{Title of the first subchapter of the first chapter}

\section{Title of the second subchapter of the first chapter}


% vim: tw=110:fo=an1lt

\chapter{Problem Definition}

In this chapter, we present an overview of the problem area and a description of what exactly we are solving.
It should also serve as an introduction for readers unfamiliar with the topic, providing explanations of terms
that we are using throughout this work.


\section{File Transfer}

Most modern operating systems share a common idea of what constitutes a filesystem. It is a collection of {\it
files} that are referenced through a hierarchical tree-like structure of {\it directories}. You can identify
any file or directory by its {\it path}, which is a string of entry names delimited by path separator
characters\footnotemark[1]. You follow the (absolute) path by starting at the root directory and always
descending into the directory named by the next path element, until the last element is remaining, naming the
target entry.

We want to be able to publish one or more subtrees of the filesystem on one computer, and access this part of
the structure from another computer over a network.


\section{Protocols}


\subsection{Sessions}

The most basic "protocol" imaginable is one computer just firing data packets at another, silently hoping that
the other will reassemble and understand them. We could, in theory, use this sort of protocol for transferring
files, but it is obviously very impractical. All but the most basic protocols will first establish a session,
in order to have a context for further communication. Extent of the session can vary, from simply marking
a collection of packets as belonging together, to providing ordered delivery, to authenticating parties and
encrypting transmitted data.

Basic session functionality is usually provided by the transport layer, specifically the Transmission Control
Protocol (TCP). This allows the communicating parties to establish a two-way stream of octets, with reliable
and ordered delivery, and flow control. Unless noted otherwise, protocols described in this work are built on
top of TCP.


\subsection{Encoding}

Protocols consist of messages exchanged by the communicating parties. There are two main ways to encode the
messages: either as human-readable text, or as custom binary data.

\subsubsection{Text-based protocols}

In text-based protocols, each message is a string of words and/or numbers, nowadays usually encoded as ASCII
or UTF-8. Many such protocols also use line separator characters to delimit messages. The main advantage of
this approach is that it is readable (and writable) by humans without specialized debugging tools - for
instance, it is possible to communicate with the other party "by hand", without software implementing the
protocol. Historically, there was also the fact that two communicating systems only had to agree on text
encoding, and not about in-memory format of numbers or other data structures.

Drawbacks include overhead of text encoding, both in terms of size and processing time spent encoding and
decoding the textual representation. Another problem is that lengths of message fields are not fixed, so
parsing and framing messages is more difficult.

\subsubsection{Binary protocols}

The alternative to text-based messages is using a format similar to in-memory encoding of data. Each message
can be described by, and parsed as, a C struct\footnotemark[2]. This results in shorter messages and simpler
encoding. If the in-memory encoding is the same as in the protocol, encoding and decoding can even be skipped
completely. Parsing and framing is trivial - every kind of message can either have a fixed number of
fixed-size fields, or it can have a fixed-size part describing lengths of all variable-size fields.

Even though today's worldwide computing environment is homogenous in many regards, care must be taken when
communicating between systems that differ in endianness. The protocol either has to prescribe the byte order,
or it must provide means for the parties to negotiate.

Unlike text-based protocols, debugging is not straightforward, humans require specialized tools that convert
the binary messages to understandable text.


\subsection{Roles}

We recognize two basic roles for communicating parties: {\it client} and {\it server}. Client is the party
that sends requests - specifies what is to be done - and server's task is to fulfil the requests and possibly
report on results.

Most protocols are modelled as client-server, with roles of the parties fixed in advance. Typically, the
server is listening for incoming connections and clients connect at their convenience. In a peer-to-peer
model, after establishing the connection, parties can act as both clients and servers at the same time,
sending requests to each other.

In this work, we use the following convention: a {\it server} will be the computer that contains a "master
copy" of a filesystem, and publishes a subtree of this filesystem on the network. A {\it client} is the
computer that wants to manipulate this published tree using the protocol. 


%%%

\footnotetext[1]{On UNIX-like systems, the path separator is a forward-slash, '/'. Windows recognizes both
forward-slash and a backslash, but backslash is used more commonly. Nevertheless, we will be using
forward-slash as a path separator everywhere.}

\footnotetext[2]{This does not mean that you can create a message by dumping the memory representation of
a corresponding C struct directly. On modern systems, fields of a struct will be memory-aligned and will
contain some amount of padding. This can differ between systems, compilers, and can even depend on particular
compiler flags.}

\chapter{Related works}

This chapter examines several popular existing protocols and general characteristics
of their implementations.

\section{FTP}

FTP, the File Transfer Protocol, is one of the oldest network protocols for file transfer, and probably the
oldest that is still widely used. It was first introduced in~1973 as RFC~454~\cite{rfc454} and the most recent
accepted standard is RFC~959~\cite{rfc959} from~1985. Several proposed extensions were published after that,
mostly describing how the standard has evolved in practical use.

With FTP, the client establishes a text-based, line-oriented control connection to the server. Through this
the user can configure session parameters, traverse and list directories and manipulate directory entries,
similar to an interactive shell. Many commands are designed for human use, e.g. the HELP command, or the LIST
command that lists contents of the current directory in a platform-dependent textual format. Actual file
transfers run on separate data connections which are negotiated on the control connection. This means that in
order to support FTP transparently, firewalls and NATs must understand the FTP protocol.

Users authenticate themselves with user name, password and optionally an account name. The protocol specifies
no encryption, so all information is transferred in plain text. RFC~4217~\cite{rfc4217} specifies a variant of
FTP that uses TLS for encryption. However, this makes FTP traffic opaque, so special care must be taken when
firewalls or NATs are involved.

\section{HTTP}

Mentioned mostly for completeness, HTTP (Hyper-Text Transfer Protocol) is at its core a protocol for
retrieving files from remote servers. Although it has limited capabilities for creating, updating and deleting
content, those are usually used in application-specific manner and not for general file manipulation. It
doesn't support directory manipulation at all. Current HTTP specification is described by
RFC~2616~\cite{rfc2616} from~1999.

HTTP is text-based. Each client request specifies a method (one of GET, PUT, POST or DELETE) and an URL on the
first line, followed by any number of headers, one per line.  Headers can specify content metadata,
instructions for proxies, client expectations etc. The headers section can be followed by a body containing
arbitrary data. Similarly, server response starts with a numeric response code, followed by headers and data
payload.

The protocol is fully stateless and has no concept of sessions beyond single request/response exchanges,
although it allows reusing a single TCP connection as an optimization.  RFC~2069~\cite{rfc2069} specifies a
simplistic authentication model, based on username/password combinations. RFC~2818~\cite{rfc2818} describes
how to use the protocol with TLS, which is commonly known as HTTPS.

\subsection{WebDAV}

WebDAV (Web Distributed Authoring and Versioning) protocol is an extension of HTTP, designed to provide the
World-Wide Web network with standardized authoring capabilities. It was first specified in~1999 as
RFC~2518~\cite{rfc2518}, with RFC~4918~\cite{rfc4918} from 2007 being the latest published specification.
As~of~2012, WebDAV and its extensions are actively maintained by several IETF working groups.

WebDAV specifies several new methods that add directory manipulation and locking functionality, as well as
corresponding headers and XML-based formats of data payloads. It also specifies detailed semantics for file
and directory operations. Despite its name, versioning is not part of WebDAV's core functionality - it is
specified as an extension~\cite{rfc3253}.

WebDAV does not modify the way HTTP operates, so it is still stateless with no session support. It uses HTTP's
authentication models and it can run under HTTPS encryption mode. It is also backwards-compatible with
HTTP-only proxies. Depending on the nature of the proxy, WebDAV traffic can either pass through unchanged, or
a failure is detected and doesn't endanger consistency of server state.

The whole of WebDAV's functionality is relatively complex: it concerns itself with things like authorship,
conflict resolution, locking policies, content properties and metadata etc. Also, it is designed to run on
WWW, which means that it has to handle unique URLs, many users, anonymous vs. authorized access, proxying, i
could go on. We don't really care about any of that. So while client implementations can be relatively simple
(this is even a requirement of the original WebDAV proposal), servers tend to be complicated.

WebDAV is natively supported in most modern OSes, but due to the complexity, implementations tend to be
slightly incompatible.

\section{SMB/CIFS}

Server Message Block was designed in ancient times by arcane wizards of Microsoft for use in their desktop
operating system, Windows. Unlike the previous protocols, it doesn't require TCP/IP. It can either run
directly over TCP, or on top of a custom underlying API called NetBIOS. This was created because at that time,
TCP/IP networks were not really common. NetBIOS can run directly on physical layer, on IPX/SPX or on TCP/IP.

SMB is mainly aimed at small local networks

\section{NFS}

\section{SFTP --- SSH File Transfer Protocol}

\section{Others}


% Ukázka použití některých konstrukcí LateXu (odkomentujte, chcete-li)
%%% Ukázka použití některých konstrukcí LaTeXu

\subsection{Ukázka \LaTeX{}u}
\label{ssec:ukazka}

This short subsection serves as an~example of basic \LaTeX{} constructs,
which can be useful for writing a~thesis.

Let us start with lists:

\begin{itemize}
\item The logo of Matfyz is displayed in figure~\ref{fig:mff}.
\item This is subsection~\ref{ssec:ukazka}.
\item Citing literature~\cite{lamport94}.
\end{itemize}

Different kinds of dashes:
red-black (short),
pages 16--22 (middle),
$45-44$ (minus),
and this is --- as you could have expected --- a~sentence-level dash,
which is the longest.
(Note that we have follwed \verb|a| by a~tilde instead of a~space
to avoid line breaks at that place.)

\newtheorem{theorem}{Theorem}
\newtheorem*{define}{Definition}	% Definice nečíslujeme, proto "*"

\begin{define}
A~{\sl Tree} is a connected graph with no cycles.
\end{define}

\begin{theorem}
This theorem is false.
\end{theorem}

\begin{proof}
False theorems do not have proofs.
\end{proof}

\begin{figure}
	\centering
	\includegraphics[width=30mm]{logo.eps}
	\caption{Logo of MFF UK}
	\label{fig:mff}
\end{figure}


\chapter*{Conclusion}
\addcontentsline{toc}{chapter}{Conclusion}

\section{Future work}


%%% Seznam použité literatury
%%% Seznam použité literatury je zpracován podle platných standardů. Povinnou citační
%%% normou pro bakalářskou práci je ISO 690. Jména časopisů lze uvádět zkráceně, ale jen
%%% v kodifikované podobě. Všechny použité zdroje a prameny musí být řádně citovány.

\def\bibname{Bibliography}
\begin{thebibliography}{99}
\addcontentsline{toc}{chapter}{\bibname}

\bibitem{lamport94}
  {\sc Lamport,} Leslie.
  \emph{\LaTeX: A Document Preparation System}.
  2. vydání.
  Massachusetts: Addison Wesley, 1994.
  ISBN 0-201-52983-1.

% Chapter 1

\bibitem{asn-basic}
	{\sc International Telecommunication Union}.
	ITU-T Recommendation X.680:
	\emph{Information technology -- Abstract Syntax Notation One (ASN.1): Specification of basic notation}.
	Geneva, Switzerland: ITU-T, 2008.
	ITU-T Recommendation series X, X.680.
	Available from: \url{http://www.itu.int/rec/T-REC-X.680-200811-I}

\bibitem{rfc5246}
	{\sc Dierks,} T. and E. {\sc Rescorla}.
	\emph{The Transport Layer Security (TLS) Protocol Version 1.2}
	[online].
	RFC 5246.
	Internet Society, 2008 [viewed 4 December 2013].
	Available from: \url{http://tools.ietf.org/html/rfc5246}

% Chapter 2

\bibitem{rfc454}
	{\sc McKenzie,} A.
	\emph{File Transfer Protocol -- meeting announcement and a new proposed document}
	[online].
	RFC 454.
	Internet Society, 1973 [viewed 4 December 2013].
	Available from: \url{http://tools.ietf.org/html/rfc454}

\bibitem{rfc959}
	{\sc Postel,} J. and J. {\sc Reynolds}.
	\emph{File Transfer Protocol}
	[online].
	RFC 959.
	Internet Society, 1985 [viewed 4 December 2013].
	Available from: \url{http://tools.ietf.org/html/rfc959}

\bibitem{rfc4217}
	{\sc Ford-Hutchinson,} P.
	\emph{Securing FTP with TLS}
	[online].
	RFC 4217.
	Internet Society, 2005 [viewed 4 December 2013].
	Available from: \url{http://tools.ietf.org/html/rfc4217}

\bibitem{rfc2616}
	{\sc Fielding,} R., et al.
	\emph{Hypertext Transfer Protocol -- HTTP/1.1}
	[online].
	RFC 2616.
	Internet Society, 1999 [viewed 4 December 2013].
	Available from: \url{http://tools.ietf.org/html/rfc2616}

\bibitem{rfc2069}
	{\sc Franks,} J., et al.
	\emph{An Extension to HTTP : Digest Access Authentication}
	[online].
	RFC 2069.
	Internet Society, 1997 [viewed 4 December 2013].
	Available from: \url{http://tools.ietf.org/html/rfc2069}

\bibitem{rfc2818}
	{\sc Rescorla,} E.
	\emph{HTTP Over TLS}
	[online].
	RFC 2818.
	Internet Society, 2000 [viewed 4 December 2013].
	Available from: \url{http://tools.ietf.org/html/rfc2818}

\bibitem{rfc2518}
	{\sc Goland,} Y., et al.
	\emph{HTTP Extensions for Distributed Authoring -- WEBDAV}
	[online].
	RFC 2518.
	Internet Society 1999 [viewed 4 December 2013].
	Available from: \url{http://tools.ietf.org/html/rfc2518}

\bibitem{rfc4918}
	{\sc Internet Engineering Task Force}.
	RFC 4918:
	\emph{HTTP Extensions for Web Distributed Authoring and Versioning (WebDAV)}
	[online].
	Edited by L. Dusseault.
	June 2007 [viewed 4 December 2013].
	Available from: \url{http://tools.ietf.org/html/rfc4918}

\bibitem{rfc3253}
	{\sc Clemm,} G., et al.
	\emph{Versioning Extensions to WebDAV (Web Distributed Authoring and Versioning)}
	[online].
	RFC 3253.
	Internet Society, 2002 [viewed 4 December 2013].
	Available from: \url{http://tools.ietf.org/html/rfc3253}

\bibitem{mssmb2}
	{\sc Microsoft Corporation}.
	[MS-SMB2]:
	\emph{Server Message Block (SMB) Protocol Versions 2 and 3}
	[online].
	Revision 44.0.
	Microsoft Corporation, 25 October 2013 [viewed 4 December 2013]

\bibitem{nfs85}
	{\sc Sandberg,} Russel et al.
	{\emph Design and Implementation of the Sun Network Filesystem}.
	USENIX, 1985.
	Available from: \url{http://pdos.csail.mit.edu/6.824-2004/papers/sandberg-nfs.pdf}

\bibitem{rfc1094}
	{\sc Nowicki,} B.
	\emph{NFS: Network File System Protocol specification}
	[online].
	RFC 1094.
	Internet Society, 1989 [viewed 4 December 2013].
	Available from: \url{http://tools.ietf.org/html/rfc1094}

\bibitem{rfc1813}
	{\sc Callaghan,} B., {\sc Pawlowski,} B., and P. {\sc Staubach}.
	\emph{NFS Version 3 Protocol Specification}
	[online].
	RFC 1813.
	Internet Society, 1995 [viewed 4 December 2013].
	Available from: \url{http://tools.ietf.org/html/rfc1813}

\bibitem{rfc1057}
	{\sc Sun Microsystems}.
	\emph{RPC: Remote Procedure Call Protocol specification: Version 2}
	[online].
	RFC 1057.
	Edited by .
	Internet Society, 1988 [viewed 4 December 2013].
	Available from: \url{http://tools.ietf.org/html/rfc1057}

\bibitem{rfc1832}
	{\sc Srinivasan,} R.
	\emph{XDR: External Data Representation Standard}
	[online].
	RFC 1832.
	Internet Society, 1995 [viewed 4 December 2013].
	Available from: \url{http://tools.ietf.org/html/rfc1832}

\bibitem{juszczak89}
	{\sc Juszczak,} Chet.
	Improving the Performance and Correctness of an NFS Server.
	In: {\emph Proceedings of the Winter 1989 USENIX Technical Conference}.
	USENIX, 1989, pp. 53-63.

\bibitem{rfc5661}
	{\sc Internet Engineering Task Force}.
	RFC 5661:
	\emph{Network File System (NFS) Version 4 Minor Version 1 Protocol}
	[online].
	Edited by S. Shepler, M. Eisler, and D. Noveck.
	January 2010 [viewed 4 December 2013].
	Available from: \url{http://tools.ietf.org/html/rfc5661}

\bibitem{rfc2339}
	{\sc The Internet Society} and {\sc Sun Microsystems}.
	RFC 2339:
	\emph{An Agreement Between the Internet Society, the IETF, and Sun Microsystems, Inc. in the matter of NFS V.4 Protocols}
	[online].
	May 1998 [viewed 4 December 2013].
	Available from: \url{http://tools.ietf.org/html/rfc2339}

\bibitem{rfc4251}
	{\sc Ylonen,} T.
	\emph{The Secure Shell (SSH) Protocol Architecture}
	[online].
	RFC 4251.
	Edited by C. Lonvick.
	January 2006 [viewed 4 December 2013].
	Available from: \url{http://tools.ietf.org/html/rfc4251}

\bibitem{secsh-email}
	{\sc Hartman,} Sam.
	Formal consultation prior to closing the secsh working group.
	In: \emph{ietf.secsh e-mail conference}
	[online].
	30 July 2006; 19:15:10 -0400 [viewed 4 December 2013; 18:32 CET].
	Message-ID: tslvepgknpd.fsf@cz.mit.edu.
	Available from: \url{http://osdir.com/ml/ietf.secsh/2006-07/msg00010.html}

\bibitem{draft-secsh-13}
	{\sc Galbraith,} J. and O. {\sc Saarenmaa}.
	\emph{SSH File Transfer Protocol, draft-ietf-secsh-filexfer-13.txt}
	[online].
	Internet-Draft.
	10 July 2006 [viewed 4 December 2013].
	Available from: \url{http://tools.ietf.org/html/draft-ietf-secsh-filexfer-13}

% Chapter 3

\bibitem{rfc4422}
	{\sc Internet Engineering Task Force}.
	RFC 4422:
	\emph{Simple Authentication and Security Layer (SASL)}
	[online].
	Edited by A. Melnikov and K. Zeilenga.
	June 2006 [viewed 4 December 2013].
	Available from: \url{http://tools.ietf.org/html/rfc4422}

\bibitem{kelsey2002}
	{\sc Kelsey,} John.
	Compression and Information Leakage of Plaintext.
	In: \emph{Fast Software Encryption, 9th International Workshop, FSE 2002,
	Leuven, Belgium, February 4-6, 2002, Revised Papers}
	Berlin: Springer Berlin Heidelberg, 2002, pp. 263--276.
	ISBN 978-3-540-45661-2.
	Available from: \url{http://www.iacr.org/cryptodb/archive/2002/FSE/3091/3091.pdf}

\bibitem{pep383}
	{\sc v. Löwis,} Martin.
	\emph{Non-decodable Bytes in System Character Interfaces}
	[online].
	PEP 383.
	Python Software Foundation, 2012 [viewed 4 December 2013].
	Available from: \url{http://www.python.org/dev/peps/pep-0383/}

\bibitem{rfc2743}
	{\sc Linn,} J.
	\emph{Generic Security Service Application Program Interface Version 2, Update 1}
	[online].
	RFC 2743.
	January 2000 [viewed 4 December 2013].
	Available from: \url{http://tools.ietf.org/html/rfc2743}

\bibitem{xwindow}
	{\sc Scheifler,} Robert W. and James {\sc Gettys}.
	\emph{X Window System: Core and extension protocols: X version 11, releases 6 and 6.1}.
	Edited by Al Mento and Donna Converse.
	Boston: Digital Press, 1997.
	ISBN 978-1-55558-148-0

\bibitem{rfcXX}
	{\sc Internet Engineering Task Force}.
	RFC XX:
	\emph{}
	[online].
	Edited by .
	 [viewed 4 December 2013].
	Available from: \url{http://tools.ietf.org/html/rfcXX}


\end{thebibliography}


%%% Tabulky v bakalářské práci, existují-li.
\chapwithtoc{List of Tables}

%%% Použité zkratky v bakalářské práci, existují-li, včetně jejich vysvětlení.
\chapwithtoc{List of Abbreviations}

%%% Přílohy k bakalářské práci, existují-li (různé dodatky jako výpisy programů,
%%% diagramy apod.). Každá příloha musí být alespoň jednou odkazována z vlastního
%%% textu práce. Přílohy se číslují.
\chapwithtoc{Attachments}

\openright
\end{document}
